\documentclass[10pt,a4paper]{amsart}
\usepackage[margin=19mm]{geometry}
\usepackage{amsmath,amssymb,mathtools}
\usepackage[scr=rsfs,cal=euler]{mathalfa}
\usepackage{xfrac}
\usepackage[unicode,hidelinks,bookmarks=false]{hyperref}
\newcommand{\ie}{i.e.\ }
\newcommand{\cf}{cf.\ }
\newcommand{\resp}{respectively\ }
\newcommand{\Subsection}{\subsection}
\providecommand{\nobreakdash}{\nobreak\hbox{-}}
\setlength{\emergencystretch}{2em}
\multlinegap0pt
\usepackage{amssymb}
\usepackage{xcolor}
\usepackage[shortlabels]{enumitem}
\newtheorem{thm}{Theorem}
\newtheorem{lem}{Lemma}
\newcommand{\Asc}{{\textsf{Asc}}}
\newcommand{\Des}{{\textsf{Des}}}
\newcommand{\F}{{\mathcal{F}}}
\newcommand{\LT}{{\mathcal{LT}}}
\newcommand{\MF}{{\mathcal{MF}}}
\newcommand{\T}{{\mathcal{T}}}
\newcommand{\cc}{{\mathrm{c}}}
\newcommand{\even}{{\textrm{even}}}
\newcommand{\odd}{{\textrm{odd}}}
\title{On Ternary Trees and Fighting Fish}
\author{Sen-Peng Eu, Tung-Shan Fu, Yu-Ren Pan}
\date{Source: arXiv:2509.16667v3 (CC BY 4.0)}
\begin{document}
\maketitle

\noindent\emph{Source and licence.} On Ternary Trees and Fighting Fish. Version arXiv:2509.16667v3 (CC BY 4.0), distributed by arXiv under the Creative Commons Attribution 4.0 International licence, http://creativecommons.org/licenses/by/4.0/, as recorded in the arXiv licence metadata for this exact version and shown on the abstract page https://arxiv.org/abs/2509.16667v3. Full licence notice: \texttt{licenses/arxiv-2509.16667-CC-BY-4.0.txt}.\par\smallskip

\noindent\textit{Editorial scope.} Counted unit: the article's thm thm:enumeration (source0.tex lines 1270--1280). The article's complete original proof of it is delivered with it (source0.tex lines 1282--1315, one \texttt{proof} environment of 34 lines). No mathematics is written, repaired or re-derived: the statement and the proof are the article's own lines, in the article's own notation, and no retained line is substituted or re-flowed.  Retained uncounted as the source-local context the proof needs: lines 282--284; lines 298--300; lines 848--850; lines 869--870; lines 878--885.  Nothing was omitted from any retained span, so the package carries no omission record.  No retained line contains a citation, so the wrapper carries no bibliography and no bibliography entry.  No retained line contains a figure environment, an image-inclusion command or drawing code, so no image is delivered and the package declares no asset dependency.  Editorial formatting only, in addition: an AMS \texttt{amsart} preamble that restates the article's own declarations and macros used by the retained lines, namely the environments \texttt{thm}, \texttt{lem} on separate counters, together with the standard packages those lines need.  The retained environments print in this document as Theorem 1, Lemma 1 in order of appearance, whereas the article prints the same items with the numbering of its own full document.  The complete native source of the article as distributed in the arXiv e-print for this exact version is bundled unchanged as \texttt{sources/185392/source0.tex}.
\begin{thm} \label{thm:TT-to-marked-fish} 
For all $n\ge 1$, there is a bijection $\phi: T\mapsto (F,Q)$ of $\T_n$ onto $\MF_n$ such that a ternary tree $T$ having $\ell$ nodes with odd abscissas and $n-\ell$ nodes with even abscissas is carried to an ordered pair $(F,Q)$, where the fighting fish $F$ has $\ell+1$ descending strips and $n-\ell$ ascending strips.
\end{thm}

\begin{thm} \label{thm:LTT-to-fish-bijection}
For all $n\ge 1$, there is a bijection $\psi: T\mapsto F$ of $\LT_n$ onto $\F_n$ such that a left ternary tree $T$ having $\ell$ nodes with odd abscissas, $n-\ell$ nodes with even abscissas and $k$ nodes with abscissa 0 is carried to a fighting fish $F$ having $\ell+1$ descending strips, $n-\ell$ ascending strips and a jaw with $k$ cells.
\end{thm}

\begin{lem} \label{lem:the-map-phi}
For all $n\ge 1$, there is a bijection $\phi: T\mapsto (F,Q)$ of $\T_n$ onto $\MF_n$ such that the topmost stem cell of the descending strip $Q$ of the fighting fish $F$ is the cell associated to the root of the ternary tree $T$. 
\end{lem}

\begin{lem} \label{lem:v-labeling-rule} For any ternary tree $T\in\T_n$ and any node $u\in T$, under the bijection $\chi_E:\T_n\rightarrow\T^E_n$ the abscissa of $u$ is even if $\lambda(u)\in\{E,W\}$, and odd if $\lambda(u)\in\{N,S\}$.
\end{lem}

\begin{lem}  \label{lem:unique-stem-cell-in-strips}
For any tree $T\in\T_n$, let $t$ be the root of $T$. Under the bijection $\phi$ in Lemma \ref{lem:the-map-phi}, let $(F,\Des(t))=\phi(T)\in\MF_n$. Then the following properties hold.
\begin{enumerate}
\item For any descending strip $Q\neq \Des(t)$ of $F$, there is a node $u$ with label $\lambda(u)\in\{N,S\}$ in $T$ such that the strip $Q$ is created by $u$ in the construction of $F$, and the cell $u$ is the unique stem cell with label $N$ or $S$ in the strip $Q$.
\item For any ascending strip $R$ of $F$, there is a node $u$ with label $\lambda(u)\in\{E,W\}$ in $T$ such that the strip $R$ is created by $u$ in the construction of $F$, and the cell $u$ is the unique stem cell with label $E$ or $W$ in the strip $R$.
\item If the tree $T$ contains $i$ nodes with odd abscissas and $j$ nodes with even abscissas then the fighting fish $F$ contains $i+1$ descending strips and $j$ ascending strips.
\end{enumerate}
\end{lem}

\begin{thm} \label{thm:enumeration}
For $n\ge 1$, there exist bijections
\begin{align}
\LT_n^{\odd} & \rightarrow \T_n\setminus \LT_n,  \label{eqn:LT_odd}\\
\LT_n^{\even}& \rightarrow \T_n. \label{eqn:LT_even}
\end{align}
Hence we have 
\begin{equation*}
|\LT_n| =\frac{2}{n+1} |\T_n|.
\end{equation*}
\end{thm}

\begin{proof}
For $0\le \ell\le n-1$, we shall establish a bijection
\begin{equation} \label{eqn:LT^odd-bijection}
\{(T,u)\mid u\in T\in\LT_{n,\ell}\mbox{ and $\alpha(u)$ is odd}\}\rightarrow\T_{n,\ell}\setminus\LT_{n,\ell}.
\end{equation}

Given a left ternary tree $T\in\LT_{n,\ell}$, let $t$ be the root of $T$ with $\lambda(t)=E$, and let $u_1,\dots,u_{\ell}$ be the nodes with odd abscissas of $T$.  By Lemma \ref{lem:v-labeling-rule}, we have $\lambda(u_i)\in\{N,S\}$ for each $i$. By the map $\psi$ in Theorem \ref{thm:LTT-to-fish-bijection}, let $F=\psi(T)\in\F_n$.  Then $F$ contains $\ell+1$ descending strips. By Lemma \ref{lem:unique-stem-cell-in-strips}(i), $\Des(u_1),\dots,\Des(u_{\ell})$ are the $\ell$ descending strips other than the jaw of $F$. By the proof of Theorem \ref{thm:LTT-to-fish-bijection}, the strip $\Des(t)$ is the jaw of $F$.
By the bijection in Theorem \ref{thm:TT-to-marked-fish}, let $T_j=\phi^{-1}(F,\Des(u_j))$ for $j=1,\dots, \ell$. Then $T_j\in\T_n\setminus\LT_n$ containing $\ell$ nodes with odd abscissas. Then we obtain the map $(T,u_j)\mapsto T_j\in\T_{n,\ell}\setminus\LT_{n,\ell}$, for $j=1,\dots, \ell$.

On the other hand, given a ternary tree $T\in\T_{n,\ell}\setminus\LT_{n,\ell}$, let $(F,Q)=\phi(T)\in\MF_n$. By Theorem \ref{thm:TT-to-marked-fish}, the fighting fish $F$ contains $\ell+1$ descending strips. By Lemma \ref{lem:the-map-phi}, the root cell of $F$ is the topmost stem cell of the strip $Q$. Since $T$ is not a left ternary tree, by the proof of Theorem \ref{thm:LTT-to-fish-bijection}, $Q$ is not the jaw of $F$. Let $t$ be the topmost stem cell in the jaw of $F$, and let $T_F=\phi^{-1}(F,\Des(t))\in\T_n$. Then $T_F$ is a left ternary tree containing $\ell$ nodes with odd abscissas. By Lemma \ref{lem:unique-stem-cell-in-strips}(i), there is a unique stem cell, say $u$, with label $N$ or $S$ in the strip $Q$ of $F$. By Lemma \ref{lem:v-labeling-rule}, the node $u$ is at odd abscissa in $T_F$. Thus, we have the inverse map $T\mapsto (T_F,u)$. The bijection (\ref{eqn:LT^odd-bijection}) is established. Note that 
\begin{align*}
\LT_n^{\odd} &= \bigcup_{\ell=0}^{n-1} 
\{(T,u)\mid u\in T\in\LT_{n,\ell}\mbox{ and $\alpha(u)$ is odd}\}, \\
\T_n\setminus \LT_n &= \bigcup_{\ell=0}^{n-1} \big(\T_{n,\ell} \setminus\LT_{n,\ell}\big).
\end{align*}
The bijection in (\ref{eqn:LT_odd}) is established.

For $0\le \ell\le n-1$, we shall establish a bijection
\begin{equation} \label{eqn:LT^even-bijection}
\{(T,u)\mid u\in T\in\LT_{n,\ell}\mbox{ and $\alpha(u)$ is even}\}\rightarrow\T_{n,n-\ell-1}.
\end{equation}

Given a left ternary tree $T\in\LT_{n,\ell}$, let $v_1,v_2,\dots,v_{n-\ell}$ be the nodes with even abscissas of $T$, where $v_1$ is the root. By Lemma \ref{lem:v-labeling-rule}, we have $\lambda(v_i)\in\{E,W\}$ for each $i$. By the map $\psi$ in Theorem \ref{thm:LTT-to-fish-bijection}, let $F=\psi(T)\in\F_n$. Then $F$ contains $n-\ell$ ascending strips.
For $1\le j\le n-\ell$, let $R_j=\Asc(v_j)$.
By Lemma \ref{lem:unique-stem-cell-in-strips}(ii), $R_1,\dots,R_{n-\ell}$ are the $n-\ell$ ascending strips of $F$. Consider the conjugation of $F$. Note that $R^{\cc}_1,\dots,R^{\cc}_{n-\ell}$ become the $n-\ell$ descending strips of $F^{\cc}$.  By Theorem \ref{thm:TT-to-marked-fish}, let $T_j=\phi^{-1}(F^{\cc},R^{\cc}_j)$ for $1\le j\le n-\ell$. Then each $T_j$ contains $n-\ell-1$ nodes with odd abscissas. Then we obtain the map $(T,v_j)\mapsto T_j\in\T_{n,n-\ell-1}$, for $j=1,\dots, n-\ell$.

On the other hand, given a ternary tree $T\in\T_{n,n-\ell-1}$, let $(F,Q)=\phi(T)\in\MF_n$. By Theorem \ref{thm:TT-to-marked-fish}, the fighting fish $F$ contains $n-\ell$ descending strips. In the conjugation  $F^{\cc}$, the strip $Q^{\cc}$ becomes one of the $n-\ell$ ascending strips.  By Theorem \ref{thm:LTT-to-fish-bijection}, let $T_{F^{\cc}}=\psi^{-1}(F^{\cc})\in\LT_n$. Then the left ternary tree $T_{F^{\cc}}$ contains $n-\ell$ nodes with even abscissas and hence $\ell$ nodes with odd abscissas. By Lemma \ref{lem:unique-stem-cell-in-strips}(ii), there is a unique stem cell, say $u$, with label $E$ or $W$ in the ascending strip $Q^{\cc}$ of $F^{\cc}$. By Lemma \ref{lem:v-labeling-rule}, the node $u$ is at even abscissa in $T_{F^{\cc}}$. Thus, we have the inverse map $T\mapsto (T_{F^{\cc}},u)$. The bijection (\ref{eqn:LT^even-bijection}) is established. Note that
\begin{align*}
\LT_n^{\even} &= \bigcup_{\ell=0}^{n-1} 
\{(T,u)\mid u\in T\in\LT_{n,\ell}\mbox{ and $\alpha(u)$ is even}\}, \\
\T_n &= \bigcup_{\ell=0}^{n-1} \T_{n,n-\ell-1}.
\end{align*}
The bijection in (\ref{eqn:LT_even}) is established. The result follows.
\end{proof}

\end{document}
